% D.82. Root, FGH, and PBC full independent proof/source/certificate audits PASS.
\subsection{The valuation direction and the full local Kummer matrix}

\begin{statement}{D.82}{Proved}{Every pair with two independent inertia directions away from five}
Let \(a,b\) be arbitrary completed Kummer classes for \(\mathbb Q\).
Assume that the vectors
\[
             \{(a_q,b_q)\bmod5:q\ne5\}
                         \quad\hbox{span }\mathbb F_5^2.
                                                               \tag{D.82.1}
\]
Their ramification may be infinite, and their five-valuations
\(A=a_5,B=b_5\) need not vanish. Write the exact local classes as
\[
 a=A\kappa(5)+\alpha\kappa(6),\qquad
 b=B\kappa(5)+\beta\kappa(6)
                         \quad\hbox{on }G_{\mathbb Q_5}.
                                                               \tag{D.82.2}
\]
Set
\[
 \begin{gathered}
 W=\operatorname{span}_{\mathbb F_5}
       \{(\bar A,\bar B),(\bar\alpha,\bar\beta)\},\\
 \epsilon(W)=
 \begin{cases}
  1,&W=0\text{ or }W\text{ is a line other than }
       \langle(1,0)\rangle,\langle(0,1)\rangle,\langle(1,1)\rangle,\\
  0,&\text{otherwise},
 \end{cases}\\
 n=\min\{v_5(A),v_5(\alpha)\}
   +\min\{v_5(B),v_5(\beta)\}
   +\min\{v_5(A-B),v_5(\alpha-\beta)\},
 \qquad v_5(0)=\infty .
 \end{gathered}                                               \tag{D.82.3}
\]
For the actual Alexander twist \(s=(a,b)\), the invariants are
\[
 \mathcal S_s\simeq(\mathbb Z/5)^{\epsilon(W)},\qquad
 \mathcal T_s=0,\qquad
 \mathcal K_s\simeq(\mathbb Z/5)^{\epsilon(W)}\times Z(n),
                                                               \tag{D.82.4}
\]
with the compact convention \(Z(n)\) of D.77.
In particular a two-dimensional \(W\) gives
\(\mathcal K_s=0\); a nonboundary line gives
\(\mathcal K_s=\mathbb Z/5\); a boundary line gives the free
compact group \(Z(n)\), \(n\ge1\); and \(W=0\) gives
\(\mathbb Z/5\times Z(n)\), \(n\ge3\).

The actual rational point \(t=10\) has \(\mathcal K_s=0\).
Its nonempty metabelian fibre is a torsor under
\(\prod_{\mathbb N}\mathbb Z_5\).
The general assertion is a compact-kernel calculation,
not a section-existence criterion. D.84--D.85 subsequently remove the away-five spanning
restriction when the complete residual pair is globally independent.
Globally dependent residual pairs, arbitrary infinitely ramified
first-cohomology ranks, and the first higher-derived obstruction
remain outside that completed classification.
\end{statement}

\begin{proof}
\paragraph{The Kummer cocycle of five factors through the actual image.}
Let \(I_0\) be D.78's actual untwisted operator image, and
write \(\kappa_5=\kappa(5)\), to distinguish it from the raw
weight-five Soul\'e character. We first prove that \(\kappa_5\)
is an actual continuous cocycle on \(I_0\).

If \(g\) is in the kernel of the operator representation
\(\rho_0\), comparison on \(1,t_1\) gives \(\chi(g)=1\);
then its beta series is one. D.39's exact raw-character
formula consequently gives
\(\kappa_m(g)=0\) for every odd Soul\'e weight \(m\ge3\).
Use precisely the Kummer definition
\cite[(1.1)--(1.2)]{DE_NakamuraSakugawaWojtkowiak2017}.
For each \(r\ge1\) choose
\[
                 m_r=1+4\cdot5^{r-1}.
                                                               \tag{D.82.5}
\]
For every \(a\) prime to five, \(a^{m_r-1}\equiv1\pmod{5^r}\).
Thus the defining element at this level satisfies
\[
 \prod_{\substack{1\le a<5^r\\5\nmid a}}
      (1-\zeta_{5^r}^{a})^{a^{m_r-1}}
       \equiv
 \prod_{\substack{1\le a<5^r\\5\nmid a}}
      (1-\zeta_{5^r}^{a})
       =\Phi_{5^r}(1)=5
       \pmod{\mathbb Q(\zeta_{5^r})^{\times 5^r}}.
                                                               \tag{D.82.6}
\]
An element of \(\ker\chi\) fixes both that cyclotomic field
and the roots of unity, so the Kummer root choices do not
affect the resulting equality of characters.
It follows that
\(\kappa_5(g)\equiv\kappa_{m_r}(g)=0\pmod{5^r}\) for every \(r\).
Hence \(\kappa_5\) vanishes on \(\ker\rho_0\).
The compact quotient property makes its descended cocycle
on \(I_0\) continuous. This argument takes no limit of
unverified global integral extensions of high-weight cochains:
all the raw characters used are on \(\ker\chi\).

Its image on the local cyclotomic kernel is primitive.
Over \(\mathbb Q_5(\zeta_5)\), normalize
\(v_{\zeta_5-1}(\zeta_5-1)=1\). Then \(v_{\zeta_5-1}(5)=4\).
The Kummer class of five is therefore nonzero modulo fifth
powers and is independent of the class of \(\zeta_5\), whose
valuation is zero. The unique degree-five subextension of
the cyclotomic \(\mathbb Z_5\)-extension has Kummer line
\(\langle\zeta_5\rangle\). Thus a fifth root of five is not
in that cyclotomic extension, and \(\kappa_5\bmod5\) is nonzero
on \(G_{\mathbb Q_5}\cap\ker\chi\). Its compact additive image
there is consequently all of \(\mathbb Z_5\).
D.78's local surjection onto \(I_0\) now proves the actual
quotient statement
\[
 Q=I_0/\{\chi=1,\ \kappa_5=0\}
       \simeq\mathbb Z_5(1)\rtimes\Gamma,\qquad
                  \Gamma=\mathbb Z_5^\times .
                                                               \tag{D.82.7}
\]
Indeed the cyclotomic map is onto and its kernel supplies
every translation, so every projected element can be
corrected to zero translation.

\paragraph{The exact local image and restriction diagram.}
Put \(\mathbf r=(A,B)\), \(\mathbf z=(\alpha,\beta)\), and
\(\Lambda=\mathbb Z_5[[u,v]]\), \(t_1=1+u,t_2=1+v\).
Condition (D.82.1) gives two pure inertia translations with
unit determinant away from five. The faithful actual-image
argument of D.78 therefore gives
\[
                          I_s=\mathbb Z_5^2\rtimes I_0.
                                                               \tag{D.82.8}
\]
After an integral local monomial gauge, (D.82.2) holds for
actual cocycles. D.78 proves that
\(\kappa(6)(\ker(G_{\mathbb Q_5}\to I_0))=\mathbb Z_5\).
The valuation cocycle is zero on that kernel by the first
paragraph. Hence the actual local image consists exactly of
\[
 I_5=\{(h\mathbf z+\kappa_5(g)\mathbf r,g):
                              h\in\mathbb Z_5,\ g\in I_0\}.
                                                               \tag{D.82.9}
\]
It is the semidirect product of \(H_5=\mathbb Z_5\mathbf z\)
with the graph of the descended cocycle
\(\kappa_5\mathbf r\). This graph is an actual subgroup.
It is not replaced by the zero-translation complement.

Let \(D=A_s^D\), and set
\[
 \begin{gathered}
 f_{\mathbf z}=t_1^\alpha t_2^\beta-1,\qquad
 R=\Lambda/(f_{\mathbf z}),\qquad J=(u,v)R,\\
 D_{\rm lin}=R^D,\qquad
 \lambda=(R/J)^D=\mathbb Q_5/\mathbb Z_5(-1).
 \end{gathered}                                               \tag{D.82.10}
\]
The \(I_0\)-action on \(R\) is now the actual twisted action
\[
 \rho_{\mathbf r}(g)
        =t_1^{A\kappa_5(g)}t_2^{B\kappa_5(g)}\rho_0(g).
                                                               \tag{D.82.11}
\]
Continuous coinduction for \(H_5\), including nonprimitive
and zero lines, is proved in D.81. Together with global
translation coinduction it identifies the local restriction
map with the inclusion-induced map
\[
 H^1(I_0,\lambda)\longrightarrow
             H^1(I_0,D_{\rm lin}) .
                                                               \tag{D.82.12}
\]
The coefficient action in the second group is (D.82.11).
The finite affine Chebotarev argument of D.78/D.80 thus gives
\[
 \mathcal S_s^*\lhook\joinrel\longrightarrow
 \mathcal R_s:=
 \ker\{H^1(I_0,\lambda)\to H^1(I_0,D_{\rm lin})\}.
                                                               \tag{D.82.13}
\]

\paragraph{The whole restriction kernel still has order at most five.}
Put \(w=t_1t_2-1\), \(P=uvw\), and
\(N_1=\{\chi=1,\kappa_5=0\}\subset I_0\).
For every element of \(N_1\), (D.82.11) is multiplication
by a beta series congruent to one modulo \(P\).
There is an actual local element \(\delta\in N_1\) with
\(\mathcal B_\delta-1=P E\), \(E\in\Lambda^\times\).
For completeness, on the local cyclotomic kernel the pair
\((\kappa_5,\kappa_3)\) has compact additive image. Conjugation
by an actual element with cyclotomic value
\(\omega\), \(\omega^2=-1\), acts by
\(\operatorname{diag}(\omega,\omega^3)\).
The integral projector
\((T-\omega)/(\omega^3-\omega)\) kills the first coordinate
and preserves a primitive second coordinate. D.77 proves
that the latter occurs locally. The resulting actual element
has the asserted unit multiplier by the same exact three
boundary restrictions and cubic coefficient as D.77.

Consequently
\[
 D_{\rm lin}^{I_0}=((R/PR)^D)^Q,\qquad
 (D_{\rm lin}/\lambda)^{I_0}=((J/PJ)^D)^Q.
                                                               \tag{D.82.14}
\]
The quotient is \(Q\), not just \(\Gamma\).
Its translation element \(h\in\mathbb Z_5(1)\) acts on
the compact quotients by multiplication by
\((t_1^A t_2^B)^h\). Its beta factor disappears modulo
\(P\), or modulo \(PJ\), respectively.

The compact connecting module \(C=PR/PJ\) and its exact
four-term sequence are computed in D.81 without a
primitivity assumption. It is \(\mathbb Z_5(5)\), or
\((\mathbb Z_5/\eta)(5)\) at an exact nonzero boundary of the unit vector.
The additional translation action is trivial on \(C\),
because its augmentation is one.
Also
\[
                      H^0(Q,\lambda)=H^1(Q,\lambda)=0.
                                                               \tag{D.82.15}
\]
To verify the first-cohomology assertion, use
\(1\to\mathbb Z_5(1)\to Q\to\Gamma\to1\).
The subgroup acts trivially on \(\lambda\).
Its \(H^1\) is the continuous homomorphism module
\(\mathbb Q_5/\mathbb Z_5(-2)\), since conjugation on the
source has weight one and the target has weight minus one.
Its \(\Gamma\)-invariants vanish, as do
\(H^0(\Gamma,\lambda)\) and \(H^1(\Gamma,\lambda)\);
the order-four subgroup has a unit difference from one
on each relevant character. The low-degree
Hochschild--Serre sequence proves (D.82.15).

Dualizing D.81's four-term sequence and taking \(Q\)-invariants
therefore injects \(\mathcal R_s\) into
\((C^D)^Q\), a group of order at most five.
This is a bound on the entire kernel, including possible
higher-order classes.

\paragraph{An exact order-five witness if that kernel is nonzero.}
Suppose \(\mathcal R_s\ne0\). In the coefficient long exact
sequence, the subgroup \(D_{\rm lin}^{I_0}=D^{I_5}\) maps
injectively into \((D_{\rm lin}/\lambda)^{I_0}\), since
\(\lambda^{I_0}=0\). D.77 computes \(D^{I_5}\) as a direct
sum of copies of \(\mathbb Q_5/\mathbb Z_5\), so it is
divisible. Its quotient \(\mathcal R_s\) has order five.
Injectivity of divisible abelian groups, or subtraction
of a divisible lift of five times a representative, gives
an order-five representative of its nonzero quotient class.

By (D.82.14), that representative is a continuous
\(\mathbb F_5\)-valued Cartier functional \(\ell\) on \(J/PJ\).
The injection into \(C^D\) shows that \(\ell(P)\ne0\).
Pull it back to the augmentation ideal
\(\mathfrak a=(u,v)\Lambda\). It has the following exact properties:
\[
 \begin{gathered}
 \ell(f_{\mathbf z}\Lambda)=0,\qquad
 \ell(P\mathfrak a)=0,\qquad \ell(P)\ne0,\\
 \ell(f_{\mathbf r}\mathfrak a)=0,\qquad
 f_{\mathbf r}=t_1^A t_2^B-1,\\
 \ell(\rho_0(k)j)=2\ell(j)\quad(j\in\mathfrak a),
       \qquad \chi(k)=2,\quad\kappa_5(k)=0.
 \end{gathered}                                               \tag{D.82.16}
\]
Such an actual \(k\) exists by (D.82.7).
For the fourth identity, use the actual quotient translation
with \(\chi=1,\kappa_5=1\). Its beta factor minus one lies
in \(P\Lambda\), so its effect on the ideal \(J\) is zero
after applying \(\ell\). This gives precisely
\(\ell(f_{\mathbf r}J)=0\).
All evaluations take place modulo five; no division by
a local depth is involved.

\paragraph{The primitive boundary test from D.79.}
We recall its stronger finite-functional form. If
\(\mathbf v\in\mathbb Z_5^2\) is primitive and its reduction
lies on a boundary line, then every order-five
Cartier functional on
\[
 J_{\mathbf v}/P J_{\mathbf v},\qquad
 J_{\mathbf v}=\ker\{\Lambda/(t^{\mathbf v}-1)\to\mathbb Z_5\},
 \]
which is fixed by an actual untwisted \(k\) with \(\chi(k)=2\),
annihilates \(P\).
Here is the exact step of that proof, so the assertion does
not rely on a rational rank computation. The actual boundary
permutations and integral gauges of D.79 preserve the
augmentation ideal and the product boundary ideal.
They reduce to a primitive line with one exponent of
valuation \(d\ge1\) and another a unit. If the line is an
exact boundary, \(P=0\) in its quotient. Otherwise that
quotient is \(\mathbb Z_5[[z]]\) and
\[
 P=Q E,\quad E\in\mathbb Z_5[[z]]^\times,\qquad
 Q=z^2((1+z)^{5^d}-1).
 \]
For the actual power \(g=k^{4\cdot5^{d-1}}\), D.79.20 proves
\[
 (\chi(g)^{-1}\rho_0(g)-1)z^3
          \equiv2z^{5^d+2}\pmod{5,zQ}.
                                                               \tag{D.82.17}
\]
The Cartier-invariant functional annihilates its left side,
and \(Q\equiv z^{5^d+2}\pmod5\). It thus annihilates \(Q\),
and then \(P\), because it annihilates \(Qz\).
This holds for any actual choice of \(k\); the exact boundary
formula and two Taylor coefficients used in D.79 determine
the required operator independently of unknown higher coefficients.

\paragraph{Exclusion for rank two and for a boundary line.}
If \(W=\mathbb F_5^2\), the unit vector \(\mathbf z\)
is primitive. Modulo five its quotient ring is
\(\mathbb F_5[[z]]\), with augmentation ideal \((z)\).
The independent valuation vector makes \(f_{\mathbf r}\)
have nonzero linear coefficient in this ring. Hence
\(f_{\mathbf r}J=(z^2)\), whereas \(P\in J^3\).
The fourth identity of (D.82.16) forces \(\ell(P)=0\),
a contradiction.

Now suppose that \(W\) is a nonzero boundary line.
If \(\mathbf z\) is primitive, the preceding primitive
boundary test applies directly to \(\ell\), again giving
a contradiction. Otherwise \(\mathbf r\) is primitive.
The fourth identity of (D.82.16) makes the restriction of
\(\ell\) to \(f_{\mathbf r}\Lambda\) factor through
\(f_{\mathbf r}\Lambda/f_{\mathbf r}\mathfrak a\).
On this compact cyclic module, \(\rho_0(k)\) has character
\(\chi(k)^3\): its multiplier has augmentation \(\chi^2\),
and \(f_{\mathbf r}(t^\chi)/f_{\mathbf r}(t)\) has
augmentation \(\chi\). Cartier invariance thus gives
\[
                  (2^3-2)\ell(f_{\mathbf r})=0.
 \]
The coefficient six is a five-adic unit. Consequently
\(\ell\) vanishes on \(f_{\mathbf r}\Lambda\), not only
on \(f_{\mathbf r}\mathfrak a\). It descends to the
primitive \(\mathbf r\)-line quotient, where (D.82.17)
forces \(\ell(P)=0\). This contradiction proves
\[
           \mathcal S_s=0
       \quad\text{if \(W\) has rank two or is a boundary line.}
                                                               \tag{D.82.18}
\]

\paragraph{The two nonzero cases have one common finite-level class.}
Let \(c=48h\bmod5\) be D.81's nonzero class with
\(\lambda[5]\)-coefficients. The global actual image
(D.82.8) makes its image in \(H^1(G_{\mathbb Q},D)\)
nonzero, by the same coinduction and inflation argument.
Use \(M=(A_s/(5,(u,v)^4))^D\).
If \(W=0\), every local translation vanishes on \(M\),
and the single functional in D.81.18 is a primitive
for the entire local class at five.

If \(W\) is a nonboundary line, choose an integral lift
\(\mathbf e=(e_1,e_2)\) with
\(e_1e_2(e_1-e_2)\in\mathbb Z_5^\times\).
All local translation vectors modulo five lie in
\(\mathbb F_5\bar{\mathbf e}\).
They therefore disappear on the actual stable quotient
\[
 \Lambda/(5,(u,v)^4,t_1^{e_1}t_2^{e_2}-1)
       \simeq\mathbb F_5[z]/(z^4),\quad
             t_1=(1+z)^{e_2},\quad t_2=(1+z)^{-e_1}.
                                                               \tag{D.82.19}
\]
Its Cartier dual is a submodule of \(M\).
The image of \(P\) has nonzero cubic coefficient
\(e_2(-e_1)(e_2-e_1)\). The unit-line cubic calculation
of D.78 consequently supplies a primitive for \(c\)
in this four-dimensional dual, for the entire local action.
One can see the existence directly: prescribe its top
functional to detect \(P\) with coefficient \(-48\);
the three lower Cartier eigenvalues of an actual
residual \((\chi,c)=(2,0)\) are \(2^{-1},2^{-2},2^{-3}\),
all with unit differences from one. There is a unique
correction in those lower degrees making it fixed by
that element. The residual \((1,1)\) element has exactly
the required constant defect. These two actual residual
operators generate the group with parameters
\((\chi,c)\in\mathbb F_5^\times\ltimes\mathbb F_5(3)\),
so the identity holds for every local element.
The accompanying certificate verifies all three
nonboundary projective lines and all their actual
translation, cyclotomic and cubic residual values.

The common finite-level global argument of D.80 now
applies without change. Away from five only finitely
many inertia vectors are nonzero modulo five.
The universal cubic-module Frobenius primitives handle
all remaining places uniformly; at each exceptional
place D.43 and finite local Tate duality kill the one
class at some larger finite level. Their finite maximum
works globally and preserves the just-constructed
primitive at five. Therefore \(c\) is a nonzero
finite-level everywhere locally trivial class. The
whole-kernel bound proves
\(\mathcal S_s=\mathbb Z/5\) in both remaining cases.

Condition (D.82.1) makes \(a,b,a-b\) all globally
primitive. The local completed Kummer basis is
\(\kappa(5),\kappa(6)\), so their local divisibilities
are exactly the three minima in (D.82.3).
D.77 proves the remaining compact statements and their
continuous product splitting. The four descriptions
by \(W\) follow directly from these minima.

\paragraph{An actual rational point and arbitrary local matrices.}
For \(t=10\), the actual point section has pair
\[
        a=\kappa(2)+\kappa(5),\qquad b=2\kappa(3).
                                                               \tag{D.82.20}
\]
The sign in \(1-t=-9\) vanishes in pro-five Kummer cohomology.
Inertia at two and three has determinant two, a unit.
Its valuation vector is \((1,0)\), while D.78's exact
unit logarithms give
\((\alpha,\beta)\equiv(2,3)\pmod5\).
Thus \(W\) has rank two and \(\mathcal K_s=0\).
The full point section proves nonemptiness, and D.74
computes its finite-support metabelian fibre as the
stated countable-product torsor.

More generally every integral pair of local vectors
\((A,B),(\alpha,\beta)\) occurs in the hypotheses:
add \(A\kappa(5),B\kappa(5)\), respectively, to
D.81.20 with its prescribed local unit values
\(\alpha,\beta\). The away-five inertia vectors at
seven and forty-three remain the standard basis.
This observation constructs actual finite-support
Galois modules in every local matrix stratum; it makes
no further section-existence assertion.
\end{proof}
